\section{Experimental Setup}
\label{sec:setup}

All experiments are conducted within the NS-3 O-RAN 6G framework described
in Section~\ref{sec:architecture}.
Table~\ref{tab:sim_config} summarises the primary simulation parameters.

%% Auto-generated by tools/paper/generate_paper_assets.py
\begin{table}[htbp]
  \centering
  \caption{Simulation Configuration}
  \label{tab:sim_config}
  \begin{tabular}{@{}ll@{}}
    \toprule
    \textbf{Parameter} & \textbf{Value} \\
    \midrule
    Number of gNBs & 7 \\
    Number of UEs & 20 \\
    Simulation time (s) & 300 \\
    Total handovers & 50 \\
    Total measurements & 42000 \\
    Handover algorithms & RL-Optimised / Distance-Based / RSRP-Based \\
    Channel model & NR sub-6\,GHz + THz \\
    Traffic model & Mixed (Video, IoT, AR/VR, FL) \\
    Seeds & 5 \\
    \bottomrule
  \end{tabular}
\end{table}


\subsection{Simulation Scenarios}
We evaluate three mobility scenarios:
\begin{itemize}
  \item \textbf{Static}: UEs placed at fixed random positions; baseline
        for PHY-layer characterisation.
  \item \textbf{Pedestrian} ($v \approx 1$--$5$\,m/s): representative of
        dense urban indoor/campus environments.
  \item \textbf{Vehicular} ($v \approx 10$--$30$\,m/s): highway/motorway
        handover stress test.
\end{itemize}

\subsection{Metrics}
Primary KPIs evaluated are:
\begin{itemize}
  \item Handover success rate and handover latency (ms).
  \item Per-UE downlink throughput (Mbps) and 95th-percentile latency.
  \item RSRP (dBm), RSRQ (dB), and SINR (dB).
  \item MEC service CPU/memory utilisation and end-to-end service latency.
  \item FL convergence episode and final reward.
  \item Digital-twin prediction accuracy (\%).
\end{itemize}

\subsection{Reproducibility}
Each configuration is repeated over five independent seeds.
Results files are saved to \texttt{results/data/} in CSV and JSON formats.
The asset generator (\texttt{tools/paper/generate\_paper\_assets.py}) reads
these files and produces all tables and figures in this paper.
The full simulation and paper-build pipeline can be reproduced by running:
\begin{verbatim}
  python tools/paper/generate_paper_assets.py
  cd paper && latexmk -pdf main.tex
\end{verbatim}
